\section{RLP 实验：公平比较的四个维度}

RLP 需要 rollout，因此每个训练 token 的计算成本高于 NTP。公平评估不能只比较 token 数，还要比较 FLOPs、起始 checkpoint、后续 SFT/RL、模型架构与 benchmark。课程依次展示 Qwen3-1.7B、compute-equivalent baseline 与 Nemotron-Nano-12B-V2，形成从小模型到混合架构的证据链。

\subsection{Qwen3-1.7B：实验条件先于柱状图}

第一组实验从 Qwen3-1.7B-Base final checkpoint 出发，在 general pretraining corpus 上做 RLP，不要求专门 reasoning dataset。Baseline 包括原始 base、相同 token 数的 continued pretraining（CPT），以及各自接受相同 SFT/RLVR 后的版本。只有先确认这些条件，才能解释“foundation 是否经 post-training 保留”。

\lecturefigure{slide-033.jpg}{Qwen3-1.7B RLP 实验的训练数据、baseline 与 post-training 条件}{Stanford 官方录像，00:36:31--00:37:34}

\begin{knowledgebox}{术语消化：CPT 与 RLP}
Continued Pretraining（CPT，继续预训练）仍使用标准 NTP loss，只是在 base checkpoint 上继续看更多 token；RLP 则把一部分计算用于 thought rollout 与 information-gain update。Token-matched CPT 控制数据量，FLOP-matched CPT 控制近似计算量，两者都需要。
\end{knowledgebox}

\subsection{Token-matched 结果与 post-training persistence}

在相同额外 token 数下，RLP 相对 base 与 CPT 都提升平均成绩；完成相同 post-training 后，优势仍存在。图左与图右的 benchmark/阶段不同，因此不能把柱高直接相减；教学重点是方向一致：RLP 建立的 reasoning foundation 没有被 alignment 洗掉。

\lecturefigure{slide-034.jpg}{Qwen3-1.7B 上 token-matched RLP 与 post-training 结果}{Stanford 官方录像，00:37:34--00:38:26}

\begin{knowledgebox}{读图：先分左右，再分颜色}
左侧比较 base-stage 模型，右侧比较接受相同 SFT/RLVR 的模型。绿色表示带 RLP 的路径，蓝色表示对应 baseline。需要同时看 Math、Science、Science Post、Overall，而不是只挑最高柱。图支持 benchmark-suite 内的平均相对提升，不证明所有能力域等比例增长。
\end{knowledgebox}

\subsection{FLOP matching：rollout 不是免费的}

RLP 为每个目标位置生成多条 thought，并重新计算 thought-conditioned likelihood；因此 1B RLP token 与 1B CPT token 的成本不同。课程把 170M RLP token 对应到约 6B FLOP-matched CPT token，即 CPT 看约 35 倍更多数据。即便如此，RLP 在报告设置中仍约有 14\% 平均相对优势。

\lecturefigure{slide-035.jpg}{Qwen3-1.7B 上的 FLOP-matched CPT 对照}{Stanford 官方录像，00:38:26--00:40:19}

可把粗略计算写成
\[
C_{\mathrm{RLP}}
\approx C_{\mathrm{NTP}}+G\,C_{\mathrm{rollout}}+G\,C_{\mathrm{score}},
\]
其中，$G$ 是 thought rollout 数，$C_{\mathrm{rollout}}$ 是生成 thought 的成本，$C_{\mathrm{score}}$ 是对 target token 评分的成本。真实 FLOP accounting 还依赖 KV reuse、thought length、batching、sequence packing 与 kernel 实现。

\begin{warningbox}{Token efficiency 与 wall-clock efficiency 不相同}
RLP 用更少数据 token 获得更高分，不代表训练更快或更便宜。Rollout 的串行生成、较低 hardware utilization、额外 activation 与通信都可能提高时间和能耗。部署决策必须同时报告 token、FLOPs、GPU-hours 与吞吐。
\end{warningbox}

\subsection{Nemotron-Nano-12B-V2：更大模型与 hybrid architecture}

第二组实验从约 19.8T token 的中间 checkpoint 出发，只增加 250M RLP token，并与训练到约 20T token 的 final base 比较。Nemotron-Nano-12B-V2 是 Mamba-Transformer hybrid，因此该实验同时测试模型规模、起始阶段与架构迁移。关键控制是 RLP 路径总共还少看约 200B token。

\lecturefigure{slide-036.jpg}{Nemotron 中间 checkpoint 上的有限数据 RLP 设置}{Stanford 官方录像，00:40:19--00:41:14}

\teachervoice{00:40:19--00:42:20，讲者强调 RLP 路径从 19.8T checkpoint 出发、仅使用 250M token，比较对象却是 20T final base。这个设计试图证明收益不是“多看了数据”，而是在更少总 token 下改变了学习目标。}

结果页显示 base-stage 平均相对提升约 35\%，science 领域提升尤其明显；相同 post-training 后仍保持绝对优势。图支持 RLP 跨 Transformer 与 hybrid model 的可迁移性，但模型和数据只覆盖有限实例，不能据此断言所有 architecture 都有同样 scaling slope。

\lecturefigure{slide-037.jpg}{RLP 在更大 Nemotron hybrid model 上的 base 与 post-training 结果}{Stanford 官方录像，00:41:14--00:42:20}

\begin{knowledgebox}{读图：相对 35\% 与绝对 3\% 不能混成一个数字}
左侧讲 base-stage average relative gain，右侧 post-trained bars 常以绝对分数比较。两种口径回答不同问题：前者看 foundation 改善比例，后者看最终可用模型仍领先多少。报告时必须保留 denominator 和 benchmark suite。
\end{knowledgebox}

\subsection{本章小结}

RLP 的实验可信度来自多层 baseline：token-matched 检查数据量，FLOP-matched 检查 rollout 成本，post-training-matched 检查收益是否持久，intermediate-checkpoint comparison 检查能否在更少总 token 下追过 final base。下一章把 RLP 放回 early-reasoning 文献谱系。

